\clearpage
\subsection{Résultats expérimentaux}

\paragraph{Protocole.}
Le tirage stratifié sélectionne cinq images par classe avec la graine \texttt{20261121}, soit 100 requêtes. Pour chaque requête, le script conserve le classement initial et simule cinq tours de retour. Une image est marquée pertinente si sa classe est celle de la requête. Les autres sont marquées négatives. Aucun participant n'a fourni de jugement. Les valeurs de \(K\) testées sont 5, 10, 20, 25 et 50. On analyse \(K=25\) ci-dessous pour garder une profondeur fixe. Si aucun positif n'est renvoyé, le script conserve le dernier classement pour les tours concernés.

\paragraph{Mesures.}
Pour une requête, \(H_K\) est le nombre de résultats pertinents parmi les \(K\) premiers et \(N_{\mathrm{rel}}\) le nombre d'images pertinentes dans l'index, hors requête.
\begin{equation}
  \mathrm{P@}K=\frac{H_K}{K},
  \qquad
  \mathrm{R@}K=\frac{H_K}{N_{\mathrm{rel}}}.
  \label{eq:mesures}
\end{equation}
Dans ce corpus, \(N_{\mathrm{rel}}\) vaut 499, ou 498 pour les deux classes ayant perdu un doublon. Les deux mesures sont d'abord calculées par requête puis moyennées sur les 100 requêtes. La précision mesure la proportion de bons résultats affichés. Le rappel mesure la part des résultats admissibles retrouvée dans ces \(K\) places.

\begin{table}[ht]
  \centering
  \begin{tabular}{@{}rSS@{}}
    \toprule
    Tour & {Précision@25 (\%)} & {Rappel@25 (\%)} \\
    \midrule
    0 & 33.08 & 1.66 \\
    1 & 50.52 & 2.53 \\
    2 & 53.28 & 2.67 \\
    3 & 55.32 & 2.77 \\
    4 & 55.64 & 2.79 \\
    5 & 55.52 & 2.78 \\
    \bottomrule
  \end{tabular}
  \caption{Moyennes à profondeur fixe pour 100 requêtes de VisionVault. Le tour 0 correspond aux poids uniformes.}
  \label{tab:resultats}
\end{table}

\begin{figure}[ht]
  \centering
  \begin{tikzpicture}
    \begin{axis}[
        width=0.87\linewidth,height=52mm,
        xmin=-0.12,xmax=5.12,ymin=0,ymax=100,
        xtick={0,1,2,3,4,5},ytick={0,25,50,75,100},
        xlabel={Tour},ylabel={Précision moyenne (\%)},
        axis lines=left,axis line style={draw=marine!75},
        tick label style={font=\small},label style={font=\small},
        ymajorgrids,major grid style={draw=grisclair}]
      \addplot[marine,line width=1pt,mark=*,mark size=2pt]
        table[x=iteration,y=precision,col sep=tab]{figures/precision-k25.tsv};
    \end{axis}
  \end{tikzpicture}
  \caption{Évolution de la précision à \(K=25\), tracée à partir des résultats expérimentaux de VisionVault.}
  \label{fig:precision}
\end{figure}

\clearpage
\paragraph{Interprétation.}
Le tableau~\ref{tab:resultats} et la figure~\ref{fig:precision} montrent un gain moyen de \num{22.44} points de pourcentage entre le début et le cinquième tour. Le premier tour apporte \num{17.44} points. La moyenne culmine au quatrième tour, puis baisse légèrement. Parmi les 100 requêtes, 81 progressent, 12 restent stables et 7 reculent. À \(K=25\), cinq requêtes commencent sans aucun positif et gardent leur classement dans la simulation.

\begin{figure}[ht]
  \centering
  \begin{subfigure}{0.98\linewidth}
    \centering
    \demopanel{\linewidth}{topk-avant.jpg}
    \caption{Tour 0. Une image de la même classe dans le Top-25.}
  \end{subfigure}\par\medskip
  \begin{subfigure}{0.98\linewidth}
    \centering
    \demopanel{\linewidth}{topk-apres.jpg}
    \caption{Tour 5. Vingt et une images de la même classe dans le Top-25.}
  \end{subfigure}
  \caption{Requête GHIM-10K \texttt{7\_3690.jpg}, classée par VisionVault. Chaque bande contient la requête et les cinq premiers rangs. Les nombres de correspondances en légende portent sur les 25 candidats.}
  \label{fig:retrieval}
\end{figure}

La figure~\ref{fig:retrieval} présente un cas favorable. Une autre requête, \texttt{18\_9393.jpg}, passe de 12 à 8 images pertinentes sur 25. Ces différences sont cohérentes avec la dépendance de la méthode à la sélection disponible à chaque tour. Le faible rappel s'explique aussi par la profondeur du classement. Même 25 résultats pertinents ne couvriraient qu'environ 5\,\% des quelque 500 images d'une classe. Les images déjà marquées peuvent réapparaître après reclassement. Une meilleure précision ne signifie donc pas qu'autant de nouvelles images pertinentes ont été découvertes. Enfin, une catégorie fournit une règle de comparaison reproductible, alors que la pertinence réelle peut dépendre d'un objet, d'une composition ou d'un détail de texture.

\subsection{Validation}
Le dernier essai complet sur Linux a réussi les 30 cas pytest, dont le parcours Playwright. Il a vérifié la connexion, les opérations individuelles et groupées, la hiérarchie, le recadrage et le redimensionnement, ainsi que la visualisation des cinq familles de descripteurs. La recherche initiale et deux tours de retour ont été exercés dans l'application réelle.

Les tests de l'API vérifient notamment qu'un retour porte sur le classement courant et qu'une suppression répercute les changements dans l'index. Une vérification indépendante a retrouvé les poids, les distances et les mesures des \num{3000} enregistrements expérimentaux conservés. La reconstruction à partir du code a donné les mêmes vecteurs et coefficients de normalisation. Le scénario de reconstruction a ensuite vérifié la connexion et le Top-25 dans Chromium. La procédure Windows est documentée, mais la campagne a été exécutée sur Linux.
